% answers.tex —— 2008 年·普通高等学校招生考试·天津卷·文科 · 参考答案与详细解析片段
% 由 tools/generate_projects.py 依来源总账自 2008/2008_数学真题_天津文_试题.md + 2008/2008_数学真题_天津文_解析正文.md + 2008/2008_数学真题_天津文_答案要点.md 生成，勿手改（改源 MD 后重跑）


\begin{gktip}
固定候选 PDF 仅含题面，无答案、解析和评分标准，原卷亦未标注分值；以下为非官方编者独立推导，官方答案字段为空。来源发布者、出版者、版次与官方身份均未核。
\end{gktip}

\begin{answer}
\gksolhead{1}{}

\textbf{题目}：

设集合 $U=\{x\in\mathbf N\mid0<x\leqslant8\}$，$S=\{1,2,4,5\}$，$T=\{3,5,7\}$，则 $S\cap(\complement_UT)=$（　）

（A）$\{1,2,4\}$

（B）$\{1,2,3,4,5,7\}$

（C）$\{1,2\}$

（D）$\{1,2,4,5,6,8\}$

\Needspace{4\baselineskip}
\textbf{解析}

\Needspace{5\baselineskip}
$U=\{1,2,\cdots,8\}$，$T=\{3,5,7\}$，故

\nopagebreak[4]
\[\complement_UT=\{1,2,4,6,8\},\qquad S\cap(\complement_UT)=\{1,2,4\}.\]

\gkanswer{答案：A。}

\end{answer}
\begin{answer}
\gksolhead{2}{}

\textbf{题目}：

设变量 $x$，$y$ 满足约束条件 $\begin{cases}x-y\geqslant0,\\ x+y\leqslant1,\\ x+2y\geqslant1,\end{cases}$ 则目标函数 $z=5x+y$ 的最大值为（　）

（A）$2$

（B）$3$

（C）$4$

（D）$5$

\Needspace{4\baselineskip}
\textbf{解析}

\Needspace{5\baselineskip}
可行域顶点为 $\left(\dfrac12,\dfrac12\right)$、$\left(\dfrac13,\dfrac13\right)$、$(1,0)$。计算 $z=5x+y$：

\nopagebreak[4]
\[3,\qquad 2,\qquad 5,\]

故最大值为 $5$。

\gkanswer{答案：D。}

\end{answer}
\begin{answer}
\gksolhead{3}{}

\textbf{题目}：

函数 $y=1+\sqrt x$（$0\leqslant x\leqslant4$）的反函数是（　）

（A）$y=(x-1)^2$（$1\leqslant x\leqslant3$）

（B）$y=(x-1)^2$（$0\leqslant x\leqslant4$）

（C）$y=x^2-1$（$1\leqslant x\leqslant3$）

（D）$y=x^2-1$（$0\leqslant x\leqslant4$）

\Needspace{4\baselineskip}
\textbf{解析}

\Needspace{5\baselineskip}
$x\in[0,4]$ 时 $y=1+\sqrt x\in[1,3]$，且 $\sqrt x=y-1$，故

\nopagebreak[4]
\[y=(x-1)^2,\qquad 1\leqslant x\leqslant3.\]

\gkanswer{答案：A。}

\end{answer}
\begin{answer}
\gksolhead{4}{}

\textbf{题目}：

若等差数列 $\{a_n\}$ 的前 $5$ 项和 $S_5=25$，且 $a_2=3$，则 $a_7=$（　）

（A）$12$

（B）$13$

（C）$14$

（D）$15$

\Needspace{4\baselineskip}
\textbf{解析}

\Needspace{5\baselineskip}
$S_5=5a_3=25$，故 $a_3=5$。公差

\nopagebreak[4]
\[d=a_3-a_2=2,\qquad a_7=a_3+4d=13.\]

\gkanswer{答案：B。}

\end{answer}
\begin{answer}
\gksolhead{5}{}

\textbf{题目}：

设 $a$，$b$ 是两条直线，$\alpha$，$\beta$ 是两个平面，则 $a\perp b$ 的一个充分条件是（　）

（A）$a\perp\alpha$，$b\parallel\beta$，$\alpha\perp\beta$

（B）$a\perp\alpha$，$b\perp\beta$，$\alpha\parallel\beta$

（C）$a\subset\alpha$，$b\perp\beta$，$\alpha\parallel\beta$

（D）$a\subset\alpha$，$b\parallel\beta$，$\alpha\perp\beta$

\Needspace{4\baselineskip}
\textbf{解析}

（A）$a\perp\alpha$，$b\parallel\beta$，$\alpha\perp\beta$ 时 $b$ 可在 $\beta$ 内取与 $a$ 平行的方向，假。

（B）$a\perp\alpha$，$b\perp\beta$，$\alpha\parallel\beta$ 时 $a\parallel b$，假。

（C）$b\perp\beta$，$\alpha\parallel\beta$ 故 $b\perp\alpha$；又 $a\subset\alpha$，故 $a\perp b$，真。

（D）$a\subset\alpha$，$b\parallel\beta$，$\alpha\perp\beta$ 不足，假。

\gkanswer{答案：C。}

\end{answer}
\begin{answer}
\gksolhead{6}{}

\textbf{题目}：

把函数 $y=\sin x$（$x\in\mathbf R$）的图象上所有的点向左平行移动 $\dfrac{\pi}3$ 个单位长度，再把所得图象上所有点的横坐标缩短到原来的 $\dfrac12$ 倍（纵坐标不变），得到的图象所表示的函数是（　）

（A）$y=\sin\left(2x-\dfrac{\pi}3\right)$，$x\in\mathbf R$

（B）$y=\sin\left(\dfrac x2+\dfrac{\pi}6\right)$，$x\in\mathbf R$

（C）$y=\sin\left(2x+\dfrac{\pi}3\right)$，$x\in\mathbf R$

（D）$y=\sin\left(2x+\dfrac{2\pi}3\right)$，$x\in\mathbf R$

\Needspace{4\baselineskip}
\textbf{解析}

\Needspace{5\baselineskip}
左移 $\dfrac{\pi}3$ 得 $y=\sin\left(x+\dfrac{\pi}3\right)$；横坐标缩为 $\dfrac12$ 得

\nopagebreak[4]
\[y=\sin\left(2x+\frac{\pi}3\right).\]

\gkanswer{答案：C。}

\end{answer}
\begin{answer}
\gksolhead{7}{}

\textbf{题目}：

设椭圆 $\dfrac{x^2}{m^2}+\dfrac{y^2}{n^2}=1$（$m>0$，$n>0$）的右焦点与抛物线 $y^2=8x$ 的焦点相同，离心率为 $\dfrac12$，则此椭圆的方程为（　）

（A）$\dfrac{x^2}{12}+\dfrac{y^2}{16}=1$

（B）$\dfrac{x^2}{16}+\dfrac{y^2}{12}=1$

（C）$\dfrac{x^2}{48}+\dfrac{y^2}{64}=1$

（D）$\dfrac{x^2}{64}+\dfrac{y^2}{48}=1$

\Needspace{4\baselineskip}
\textbf{解析}

\Needspace{5\baselineskip}
抛物线 $y^2=8x$ 的焦点为 $(2,0)$，故椭圆 $c=2$；$e=\dfrac12$ 给 $a=4$，

\nopagebreak[4]
\[b^2=a^2-c^2=12,\]

故 $\dfrac{x^2}{16}+\dfrac{y^2}{12}=1$。

\gkanswer{答案：B。}

\end{answer}
\begin{answer}
\gksolhead{8}{}

\textbf{题目}：

已知函数 $f(x)=\begin{cases}x+2,&x\leqslant0,\\ -x+2,&x>0,\end{cases}$ 则不等式 $f(x)\geqslant x^2$ 的解集为（　）

（A）$[-1,1]$

（B）$[-2,2]$

（C）$[-2,1]$

（D）$[-1,2]$

\Needspace{4\baselineskip}
\textbf{解析}

$x\leqslant0$ 时 $x+2\geqslant x^2$，即 $x^2-x-2\leqslant0$，得 $-1\leqslant x\leqslant0$。

$x>0$ 时 $-x+2\geqslant x^2$，即 $x^2+x-2\leqslant0$，得 $0<x\leqslant1$。

合得 $[-1,1]$。

\gkanswer{答案：A。}

\end{answer}
\begin{answer}
\gksolhead{9}{}

\textbf{题目}：

设 $a=\sin\dfrac{5\pi}7$，$b=\cos\dfrac{2\pi}7$，$c=\tan\dfrac{2\pi}7$，则（　）

（A）$a<b<c$

（B）$a<c<b$

（C）$b<c<a$

（D）$b<a<c$

\Needspace{4\baselineskip}
\textbf{解析}

\Needspace{5\baselineskip}
计算如下：

\nopagebreak[4]
\[a=\sin\frac{5\pi}7=\sin\frac{2\pi}7\approx0.782,\quad b=\cos\frac{2\pi}7\approx0.623,\quad c=\tan\frac{2\pi}7\approx1.254,\]

故 $b<a<c$。

\gkanswer{答案：D。}

\end{answer}
\begin{answer}
\gksolhead{10}{}

\textbf{题目}：

设 $a>1$，若对于任意的 $x\in[a,2a]$，都有 $y\in[a,a^2]$ 满足方程 $\log_ax+\log_ay=3$，这时 $a$ 的取值的集合为（　）

（A）$\{a\mid1<a\leqslant2\}$

（B）$\{a\mid a\geqslant2\}$

（C）$\{a\mid2\leqslant a\leqslant3\}$

（D）$\{2,3\}$

\Needspace{4\baselineskip}
\textbf{解析}

\Needspace{5\baselineskip}
$\log_ax+\log_ay=3$ 给 $y=\dfrac{a^3}x$。$x\in[a,2a]$ 时

\nopagebreak[4]
\[y\in\left[\frac{a^2}2,a^2\right].\]

要求该区间含于 $[a,a^2]$，即 $\dfrac{a^2}2\geqslant a$，故 $a\geqslant2$。

\gkanswer{答案：B。}

\end{answer}
\begin{answer}
\gksolhead{11}{}

\textbf{题目}：

一个单位共有职工 $200$ 人，其中不超过 $45$ 岁的有 $120$ 人，超过 $45$ 岁的有 $80$ 人。为了调查职工的健康状况，用分层抽样的方法从全体职工中抽取一个容量为 $25$ 的样本，应抽取超过 $45$ 岁的职工\gkblank{}人。

\Needspace{4\baselineskip}
\textbf{解析}

\Needspace{5\baselineskip}
按比例抽样：

\nopagebreak[4]
\[25\times\frac{80}{200}=10.\]

\gkanswer{答案：$10$。}

\end{answer}
\begin{answer}
\gksolhead{12}{}

\textbf{题目}：

$\left(x+\dfrac2x\right)^5$ 的二项展开式中 $x^3$ 的系数为\gkblank{}。（用数字作答）

\Needspace{4\baselineskip}
\textbf{解析}

\Needspace{5\baselineskip}
通项

\nopagebreak[4]
\[\mathrm C_5^kx^{5-k}\left(\frac2x\right)^k=\mathrm C_5^k2^kx^{5-2k}.\]

\Needspace{5\baselineskip}
令 $5-2k=3$ 得 $k=1$，系数为

\nopagebreak[4]
\[\mathrm C_5^1\times2=10.\]

\gkanswer{答案：$10$。}

\end{answer}
\begin{answer}
\gksolhead{13}{}

\textbf{题目}：

若一个球的体积为 $4\sqrt3\pi$，则它的表面积为\gkblank{}。

\Needspace{4\baselineskip}
\textbf{解析}

\Needspace{5\baselineskip}
计算如下：

\nopagebreak[4]
\[\frac43\pi R^3=4\sqrt3\pi\ \Longrightarrow\ R^3=3\sqrt3,\qquad R=\sqrt3.\]

\Needspace{5\baselineskip}
表面积

\nopagebreak[4]
\[S=4\pi R^2=12\pi.\]

\gkanswer{答案：$12\pi$。}

\end{answer}
\begin{answer}
\gksolhead{14}{}

\textbf{题目}：

已知平面向量 $\boldsymbol a=(2,4)$，$\boldsymbol b=(-1,2)$，若 $\boldsymbol c=\boldsymbol a-(\boldsymbol a\cdot\boldsymbol b)\boldsymbol b$，则 $|\boldsymbol c|=$\gkblank{}。

\Needspace{4\baselineskip}
\textbf{解析}

\Needspace{5\baselineskip}
计算如下：

\nopagebreak[4]
\[\boldsymbol a\cdot\boldsymbol b=-2+8=6,\qquad \boldsymbol c=(2,4)-6(-1,2)=(8,-8),\]

故 $|\boldsymbol c|=8\sqrt2$。

\gkanswer{答案：$8\sqrt2$。}

\end{answer}
\begin{answer}
\gksolhead{15}{}

\textbf{题目}：

已知圆 $C$ 的圆心与点 $P(-2,1)$ 关于直线 $y=x+1$ 对称。直线 $3x+4y-11=0$ 与圆 $C$ 相交于 $A$，$B$ 两点，且 $|AB|=6$，则圆 $C$ 的方程为\gkblank{}。

\Needspace{4\baselineskip}
\textbf{解析}

\Needspace{5\baselineskip}
点 $P(-2,1)$ 关于 $x-y+1=0$ 的对称点为

\nopagebreak[4]
\[\left(1-1,\ -2+1\right)=(0,-1),\]

\Needspace{5\baselineskip}
即圆心 $C(0,-1)$。圆心到直线 $3x+4y-11=0$ 的距离

\nopagebreak[4]
\[d=\frac{|0-4-11|}5=3.\]

\Needspace{5\baselineskip}
弦长 $6$ 给半弦长 $3$，故

\nopagebreak[4]
\[r^2=d^2+3^2=18,\]

圆为 $x^2+(y+1)^2=18$。

\gkanswer{答案：$x^2+(y+1)^2=18$。}

\end{answer}
\begin{answer}
\gksolhead{16}{}

\textbf{题目}：

有 $4$ 张分别标有数字 $1$，$2$，$3$，$4$ 的红色卡片和 $4$ 张分别标有数字 $1$，$2$，$3$，$4$ 的蓝色卡片，从这 $8$ 张卡片中取出 $4$ 张卡片排成一行。如果取出的 $4$ 张卡片所标的数字之和等于 $10$，则不同的排法共有\gkblank{}种。（用数字作答）

\Needspace{4\baselineskip}
\textbf{解析}

\Needspace{5\baselineskip}
所取四张卡片的数字（颜色可重复）之和为 $10$，可能的多重集只有

\nopagebreak[4]
\[\{4,4,1,1\},\qquad\{4,3,2,1\},\qquad\{3,3,2,2\}.\]

$\{4,4,1,1\}$：必须取两张 $4$ 与两张 $1$，四个不同卡片，排列 $4!=24$ 种。

$\{4,3,2,1\}$：每个数字取红、蓝之一，颜色 $2^4=16$ 种，排列 $4!$ 种，共 $384$ 种。

$\{3,3,2,2\}$：两张 $3$ 与两张 $2$，排列 $4!=24$ 种。

\Needspace{5\baselineskip}
合计

\nopagebreak[4]
\[24+384+24=432.\]

\gkanswer{答案：$432$。}

\end{answer}
\begin{answer}
\gksolhead{17}{}

\textbf{题目}：

已知函数 $f(x)=2\cos^2\omega x+2\sin\omega x\cos\omega x+1$（$x\in\mathbf R$，$\omega>0$）的最小正周期是 $\dfrac{\pi}2$。

（1）求 $\omega$ 的值；

（2）求函数 $f(x)$ 的最大值，并且求使 $f(x)$ 取得最大值的 $x$ 的集合。

\Needspace{4\baselineskip}
\textbf{解析}

\Needspace{5\baselineskip}
（1）

\nopagebreak[4]
\[f(x)=(1+\cos2\omega x)+\sin2\omega x+1=2+\sqrt2\sin\left(2\omega x+\frac{\pi}4\right).\]

最小正周期 $\dfrac{2\pi}{2\omega}=\dfrac{\pi}\omega=\dfrac{\pi}2$，故 $\omega=2$。

\Needspace{5\baselineskip}
（2）$f(x)=2+\sqrt2\sin\left(4x+\dfrac{\pi}4\right)$，最大值 $2+\sqrt2$，当

\nopagebreak[4]
\[4x+\frac{\pi}4=\frac{\pi}2+2k\pi\ \Longrightarrow\ x=\frac{\pi}{16}+\frac{k\pi}2.\]

\gkanswer{答案：（1）$\omega=2$；（2）最大值 $2+\sqrt2$，$x\in\left\{\dfrac{\pi}{16}+\dfrac{k\pi}2\ \middle|\ k\in\mathbf Z\right\}$。}

\end{answer}
\begin{answer}
\gksolhead{18}{}

\textbf{题目}：

甲、乙两个篮球运动员互不影响地在同一位置投球，命中率分别为 $\dfrac12$ 与 $p$，且乙投球 $2$ 次均未命中的概率为 $\dfrac1{16}$。

（1）求乙投球的命中率 $p$；

（2）求甲投球 $2$ 次，至少命中 $1$ 次的概率；

（3）若甲、乙两人各投球 $2$ 次，求两人共命中 $2$ 次的概率。

\Needspace{4\baselineskip}
\textbf{解析}

\Needspace{5\baselineskip}
（1）乙两次均未命中：

\nopagebreak[4]
\[(1-p)^2=\frac1{16}\ \Longrightarrow\ 1-p=\frac14,\qquad p=\frac34.\]

\Needspace{5\baselineskip}
（2）甲两次至少命中一次：

\nopagebreak[4]
\[1-\left(\frac12\right)^2=\frac34.\]

\Needspace{5\baselineskip}
（3）设甲、乙命中次数为 $X\sim B\left(2,\dfrac12\right)$，$Y\sim B\left(2,\dfrac34\right)$，

\nopagebreak[4]
\[P(X=0)=\frac14,\quad P(X=1)=\frac12,\quad P(X=2)=\frac14,\]

\nopagebreak[4]
\[P(Y=0)=\frac1{16},\quad P(Y=1)=\frac38,\quad P(Y=2)=\frac9{16}.\]

\nopagebreak[4]
\[P(X+Y=2)=\frac14\cdot\frac9{16}+\frac12\cdot\frac38+\frac14\cdot\frac1{16}
=\frac{9+12+1}{64}=\frac{11}{32}.\]

\gkanswer{答案：（1）$p=\dfrac34$；（2）$\dfrac34$；（3）$\dfrac{11}{32}$。}

\end{answer}
\begin{answer}
\gksolhead{19}{}

\textbf{题目}：

如图，在四棱锥 $P-ABCD$ 中，底面 $ABCD$ 是矩形。已知 $AB=3$，$AD=2$，$PA=2$，$PD=2\sqrt2$，$\angle PAB=60^\circ$。

（1）证明 $AD\perp$ 平面 $PAB$；

（2）求异面直线 $PC$ 与 $AD$ 所成的角的大小；

（3）求二面角 $P-BD-A$ 的大小。
\nopagebreak[4]
\begin{center}% 2008 年天津卷文科第 19 题：四棱锥 P-ABCD，底面 ABCD 是矩形，
%   AB=3，AD=2，PA=2，PD=2√2，∠PAB=60°。
% 来源 PDF：gaokaomath-src/普通高考/2008/2008天津文.pdf
% SHA-256：4c8f1c998f95d805901f2e0d93fdbb5901ce4bc67fd7c3ca712083325837bb3d
% 原页：2008天津文.pdf 第 1 页（A3 三栏）第 3 栏，书页码 489。
% 重绘依据：原图底面为平行四边形 ABCD，A 为后侧顶点；P—B、P—C、P—D、
%   B—C、C—D 实线，P—A、A—B、A—D、B—D 虚线，A 处有 PA⊥AD 的直角标记。
%   坐标按原图 300dpi 像素位置等比换算（B(540,510)、C(825,515)、D(975,370)、
%   A(685,370)、P(640,170)）。
\begin{tikzpicture}[scale=1.0, semithick,
  every node/.style={inner sep=1.5pt}]
  \coordinate (B) at (0.00,0.00);
  \coordinate (C) at (2.85,0.00);
  \coordinate (D) at (4.35,1.45);
  \coordinate (A) at (1.45,1.45);
  \coordinate (P) at (1.00,3.45);

  \draw (P) -- (B);
  \draw (P) -- (C);
  \draw (P) -- (D);
  \draw (B) -- (C);
  \draw (C) -- (D);
  \draw[dashed] (P) -- (A);
  \draw[dashed] (A) -- (B);
  \draw[dashed] (A) -- (D);
  \draw[dashed] (B) -- (D);

  % A 处 PA⊥AD 的直角标记
  \coordinate (A1) at ($(A)!0.26cm!(P)$);
  \coordinate (A2) at ($(A)!0.26cm!(D)$);
  \draw (A1) -- ($(A1)+(A2)-(A)$) -- (A2);

  \node[below left=1pt] at (A) {$A$};
  \node[below left] at (B) {$B$};
  \node[below right] at (C) {$C$};
  \node[right] at (D) {$D$};
  \node[above] at (P) {$P$};
\end{tikzpicture}
\end{center}
\nopagebreak[4]
\Needspace{4\baselineskip}
\textbf{解析}

\Needspace{5\baselineskip}
（1）矩形给 $AD\perp AB$。又

\nopagebreak[4]
\[PA^2+AD^2=4+4=8=PD^2,\]

\Needspace{5\baselineskip}
故 $PA\perp AD$。$AD$ 垂直于平面 $PAB$ 内两条相交直线 $AB$，$PA$，故

\nopagebreak[4]
\[AD\perp\text{\text{平面} }PAB.\]

\Needspace{5\baselineskip}
（2）以 $A$ 为原点，$AB$ 为 $x$ 轴，$AD$ 为 $y$ 轴。由（1）平面 $PAB$ 为 $y=0$；$PA=2$，$\angle PAB=60^\circ$，故

\nopagebreak[4]
\[P(1,0,\sqrt3),\quad B(3,0,0),\quad C(3,2,0),\quad D(0,2,0).\]

\nopagebreak[4]
\[\overrightarrow{PC}=(2,2,-\sqrt3),\qquad \overrightarrow{AD}=(0,2,0).\]

\nopagebreak[4]
\[\cos\theta=\frac{\left|\overrightarrow{PC}\cdot\overrightarrow{AD}\right|}{\left|\overrightarrow{PC}\right|\left|\overrightarrow{AD}\right|}
=\frac4{2\sqrt{11}}=\frac2{\sqrt{11}},\]

故所成角为 $\arccos\dfrac2{\sqrt{11}}$。

\Needspace{5\baselineskip}
（3）平面 $ABD$ 的法向量 $\boldsymbol n_1=(0,0,1)$。平面 $PBD$ 内

\nopagebreak[4]
\[\overrightarrow{BD}=(-3,2,0),\qquad \overrightarrow{BP}=(-2,0,\sqrt3),\]

\nopagebreak[4]
\[\boldsymbol n_2=\overrightarrow{BD}\times\overrightarrow{BP}=\left(2\sqrt3,3\sqrt3,4\right).\]

\nopagebreak[4]
\[\cos\varphi=\frac{\left|\boldsymbol n_1\cdot\boldsymbol n_2\right|}{\left|\boldsymbol n_1\right|\left|\boldsymbol n_2\right|}
=\frac4{\sqrt{55}},\]

故二面角为 $\arccos\dfrac4{\sqrt{55}}$。

\gkanswer{答案：（1）证明见解析；（2）$\arccos\dfrac2{\sqrt{11}}$；（3）$\arccos\dfrac4{\sqrt{55}}$。}

\end{answer}
\begin{answer}
\gksolhead{20}{}

\textbf{题目}：

已知数列 $\{a_n\}$ 中，$a_1=1$，$a_2=2$，且 $a_{n+1}=(1+q)a_n-qa_{n-1}$（$n\geqslant2$，$q\neq0$）。

（1）设 $b_n=a_{n+1}-a_n$（$n\in\mathbf N^*$），证明 $\{b_n\}$ 是等比数列；

（2）求数列 $\{a_n\}$ 的通项公式；

（3）若 $a_3$ 是 $a_6$ 与 $a_9$ 的等差中项，求 $q$ 的值，并证明：对任意的 $n\in\mathbf N^*$，$a_n$ 是 $a_{n+3}$ 与 $a_{n+6}$ 的等差中项。

\Needspace{4\baselineskip}
\textbf{解析}

\Needspace{5\baselineskip}
（1）由 $a_{n+1}=(1+q)a_n-qa_{n-1}$ 得

\nopagebreak[4]
\[b_n=a_{n+1}-a_n=q(a_n-a_{n-1})=qb_{n-1}.\]

又 $b_1=a_2-a_1=1$，故 $\{b_n\}$ 是首项 $1$、公比 $q$ 的等比数列。

\Needspace{5\baselineskip}
（2）

\nopagebreak[4]
\[a_n=a_1+\sum_{k=1}^{n-1}b_k=1+\sum_{k=1}^{n-1}q^{k-1}.\]

\Needspace{5\baselineskip}
$q\neq1$ 时

\nopagebreak[4]
\[a_n=1+\frac{q^{n-1}-1}{q-1};\]

$q=1$ 时 $a_n=n$。

\Needspace{5\baselineskip}
（3）$a_3=2+q$，$a_6=q^4+q^3+q^2+q+2$，$a_9=q^7+q^6+q^5+q^4+q^3+q^2+q+2$。由 $2a_3=a_6+a_9$：

\nopagebreak[4]
\[4+2q=q^7+q^6+q^5+2q^4+2q^3+2q^2+2q+4,\]

\nopagebreak[4]
\[q^2\left(q^5+q^4+q^3+2q^2+2q+2\right)=0.\]

\Needspace{5\baselineskip}
$q\neq0$，且

\nopagebreak[4]
\[q^5+q^4+q^3+2q^2+2q+2=(q^2+q+1)(q^3+2),\]

$q^2+q+1>0$，故 $q^3=-2$，$q=-\sqrt[3]{2}$。

\Needspace{5\baselineskip}
此时对任意 $n$，

\nopagebreak[4]
\[a_{n+3}+a_{n+6}=2a_n+q^{n-1}\left[(1+q+q^2)+(1+q+\cdots+q^5)\right].\]

\Needspace{5\baselineskip}
括号内为

\nopagebreak[4]
\[2+2q+2q^2+q^3+q^4+q^5=(q^2+q+1)(q^3+2)=0,\]

故 $a_{n+3}+a_{n+6}=2a_n$，即 $a_n$ 是 $a_{n+3}$ 与 $a_{n+6}$ 的等差中项。

\gkanswer{答案：（1）证明见解析；（2）$a_n=1+\dfrac{q^{n-1}-1}{q-1}$（$q\neq1$），$a_n=n$（$q=1$）；（3）$q=-\sqrt[3]{2}$，证明见解析。}

\end{answer}
\begin{answer}
\gksolhead{21}{}

\textbf{题目}：

设函数 $f(x)=x^4+ax^3+2x^2+b$（$x\in\mathbf R$），其中 $a$，$b\in\mathbf R$。

（1）当 $a=-\dfrac{10}3$ 时，讨论函数 $f(x)$ 的单调性；

（2）若函数 $f(x)$ 仅在 $x=0$ 处有极值，求 $a$ 的取值范围；

（3）若对于任意的 $a\in[-2,2]$，不等式 $f(x)\leqslant1$ 在 $[-1,1]$ 上恒成立，求 $b$ 的取值范围。

\Needspace{4\baselineskip}
\textbf{解析}

\Needspace{5\baselineskip}
（1）$a=-\dfrac{10}3$ 时

\nopagebreak[4]
\[f'(x)=4x^3-10x^2+4x=2x(2x-1)(x-2).\]

符号：$x<0$ 负；$0<x<\dfrac12$ 正；$\dfrac12<x<2$ 负；$x>2$ 正。故减区间为 $(-\infty,0)$、$\left(\dfrac12,2\right)$，增区间为 $\left(0,\dfrac12\right)$、$(2,+\infty)$。

\Needspace{5\baselineskip}
（2）

\nopagebreak[4]
\[f'(x)=x\left(4x^2+3ax+4\right).\]

\Needspace{5\baselineskip}
$x=0$ 是极值点（$f'$ 在 $0$ 两侧变号）。要使 $x=0$ 是唯一极值点，二次式

\nopagebreak[4]
\[4x^2+3ax+4\]

\Needspace{5\baselineskip}
不能有使 $f'$ 变号的实根：判别式

\nopagebreak[4]
\[\Delta=9a^2-64\leqslant0\ \Longrightarrow\ |a|\leqslant\frac83.\]

（$\Delta=0$ 时为重根，$f'$ 不变号，仍无新增极值。）

\Needspace{5\baselineskip}
（3）需对一切 $x\in[-1,1]$、$a\in[-2,2]$：

\nopagebreak[4]
\[b\leqslant1-x^4-2x^2-ax^3.\]

\Needspace{5\baselineskip}
最坏情形为 $a x^3=2|x|^3$，故需

\nopagebreak[4]
\[b\leqslant1-t^4-2t^2-2t^3\qquad(t=|x|\in[0,1]).\]

\Needspace{5\baselineskip}
右边关于 $t$ 递减，最小值在 $t=1$：

\nopagebreak[4]
\[b\leqslant1-1-2-2=-4.\]

\gkanswer{答案：（1）减区间 $(-\infty,0)$、$\left(\dfrac12,2\right)$，增区间 $\left(0,\dfrac12\right)$、$(2,+\infty)$；（2）$|a|\leqslant\dfrac83$；（3）$b\leqslant-4$。}

\end{answer}
\begin{answer}
\gksolhead{22}{}

\textbf{题目}：

已知中心在原点的双曲线 $C$ 的一个焦点是 $F_1(-3,0)$，一条渐近线的方程是 $\sqrt5x-2y=0$。

（1）求双曲线 $C$ 的方程；

（2）若以 $k$（$k\neq0$）为斜率的直线 $l$ 与双曲线 $C$ 相交于两个不同的点 $M$，$N$，线段 $MN$ 的垂直平分线与两坐标轴围成的三角形的面积为 $\dfrac{81}2$，求 $k$ 的取值范围。

\Needspace{4\baselineskip}
\textbf{解析}

\Needspace{5\baselineskip}
（1）$c=3$，渐近线斜率 $\dfrac ba=\dfrac{\sqrt5}2$，故 $b^2=\dfrac54a^2$。由 $a^2+b^2=9$：

\nopagebreak[4]
\[\frac94a^2=9\ \Longrightarrow\ a^2=4,\quad b^2=5,\]

双曲线为 $\dfrac{x^2}4-\dfrac{y^2}5=1$。

\Needspace{5\baselineskip}
（2）设 $l:y=kx+m$。代入双曲线：

\nopagebreak[4]
\[(5-4k^2)x^2-8kmx-4m^2-20=0.\]

\Needspace{5\baselineskip}
两交点要求 $k^2\neq\dfrac54$ 且

\nopagebreak[4]
\[\Delta>0\ \Longleftrightarrow\ 5m^2+25-20k^2>0\ \Longleftrightarrow\ m^2>4k^2-5.\]

\Needspace{5\baselineskip}
设中点 $(x_0,y_0)$，由韦达定理

\nopagebreak[4]
\[x_0=\frac{4km}{5-4k^2},\qquad y_0=\frac{5m}{5-4k^2}.\]

\Needspace{5\baselineskip}
垂直平分线斜率为 $-\dfrac1k$，与两轴的交点横、纵坐标分别为

\nopagebreak[4]
\[X=x_0+ky_0=\frac{9km}{5-4k^2},\qquad Y=y_0+\frac{x_0}k=\frac{9m}{5-4k^2}.\]

\Needspace{5\baselineskip}
三角形面积

\nopagebreak[4]
\[\frac12|XY|=\frac{81|k|m^2}{2(5-4k^2)^2}=\frac{81}2,\]

\Needspace{5\baselineskip}
故

\nopagebreak[4]
\[m^2=\frac{(5-4k^2)^2}{|k|}.\]

\Needspace{5\baselineskip}
代入 $m^2>4k^2-5$。当 $k^2<\dfrac54$ 时右端为负，恒成立；当 $k^2>\dfrac54$ 时需

\nopagebreak[4]
\[(4k^2-5)^2>(4k^2-5)\sqrt{k^2}\ \Longrightarrow\ 4k^2-5>\sqrt{k^2}.\]

令 $v=|k|$：$4v^2-v-5>0$，即 $(4v-5)(v+1)>0$，故 $v>\dfrac54$。

\Needspace{5\baselineskip}
综上

\nopagebreak[4]
\[k\in\left(-\infty,-\frac54\right)\cup\left(-\frac{\sqrt5}2,0\right)\cup\left(0,\frac{\sqrt5}2\right)\cup\left(\frac54,+\infty\right).\]

\gkanswer{答案：（1）$\dfrac{x^2}4-\dfrac{y^2}5=1$；（2）$\left(-\infty,-\dfrac54\right)\cup\left(-\dfrac{\sqrt5}2,0\right)\cup\left(0,\dfrac{\sqrt5}2\right)\cup\left(\dfrac54,+\infty\right)$。}

\end{answer}